% Ecuación de Slutsky para el precio propio, versión reducida para la pizarra del vídeo
% (la tabla completa, con el bien frontera, está en el tema: tab:slutsky-propio).
% Cada fila aparece en su capa: bien normal, inferior ordinario e inferior Giffen.
% Autor: Víctor Gutiérrez Marcos
\documentclass[tikz,border=4pt]{standalone}
% ==========================================================================
% tikz-tcee.tex — Estilos comunes de los gráficos TikZ del temario
% Autor: Víctor Gutiérrez Marcos
%
% Cada gráfico es un fichero standalone en fuentes/<tema>/graficos/:
%
%   \documentclass[tikz,border=4pt]{standalone}
%   % ==========================================================================
% tikz-tcee.tex — Estilos comunes de los gráficos TikZ del temario
% Autor: Víctor Gutiérrez Marcos
%
% Cada gráfico es un fichero standalone en fuentes/<tema>/graficos/:
%
%   \documentclass[tikz,border=4pt]{standalone}
%   % ==========================================================================
% tikz-tcee.tex — Estilos comunes de los gráficos TikZ del temario
% Autor: Víctor Gutiérrez Marcos
%
% Cada gráfico es un fichero standalone en fuentes/<tema>/graficos/:
%
%   \documentclass[tikz,border=4pt]{standalone}
%   \input{../../../latex/tikz-tcee}
%   \begin{document}
%   \begin{tikzpicture}[tcee]
%     \ejes{Q}{P}{6}{5}
%     \draw[curva1] (0.5,4.5) -- (5,0.8) node[etiqueta,right] {$D$};
%     \capa{2}{\draw[curva2] ...;}   % en el vídeo aparece en el paso 2
%   \end{tikzpicture}
%   \end{document}
%
% Paleta: la de la web (morado, dorado, salvia) más tres colores de apoyo
% distinguibles en blanco y negro por el trazo.
% ==========================================================================
\usepackage[T1]{fontenc}
\usepackage{newpxtext,newpxmath}
\usepackage{xcolor}
\usepackage{tikz,pgfplots}
\pgfplotsset{compat=1.18}
\usetikzlibrary{arrows.meta,calc,positioning,patterns,decorations.pathreplacing,intersections,shapes.misc,backgrounds}

\definecolor{tceemorado}{HTML}{5F2987}
\definecolor{tceedorado}{HTML}{B8860B}
\definecolor{tceeverde}{HTML}{3C7D3A}
\definecolor{tceeazul}{HTML}{1F5F99}
\definecolor{tceerojo}{HTML}{B22222}
\definecolor{tceegris}{HTML}{767171}
\definecolor{tceesalvia}{HTML}{E2EFD9}

% Capas para el vídeo: \capa{n}{…} dibuja su contenido solo a partir del paso n.
% En el PDF, la web y Word se ve todo; generar-video.py compila el gráfico una
% vez por paso con \def\tceepaso{k} para que cada elemento aparezca al narrarlo.
\providecommand{\tceepaso}{1000}
\newcommand{\capa}[2]{\ifnum#1>\tceepaso\relax\else#2\fi}

\tikzset{
  tcee/.style={line cap=round,line join=round,font=\small,>={Stealth[length=2.2mm]}},
  eje/.style={->,thick,black},
  curva1/.style={very thick,tceemorado},
  curva2/.style={very thick,tceeazul},
  curva3/.style={very thick,tceeverde},
  curva4/.style={very thick,tceedorado},
  curvanueva/.style={very thick,tceerojo},
  desplazada/.style={very thick,dashed},
  guia/.style={thin,dashed,tceegris},
  punto/.style={circle,fill=black,inner sep=1.3pt},
  etiqueta/.style={font=\small,inner sep=2pt},
  flecha/.style={->,thick,tceerojo},
  area/.style={fill=tceemorado,fill opacity=0.18,draw=none},
  area2/.style={fill=tceedorado,fill opacity=0.22,draw=none},
}

% \ejes{etiqueta x}{etiqueta y}{largo x}{alto y}
\newcommand{\ejes}[4]{%
  \draw[eje] (0,0) -- (#3,0) node[below left=2pt and 0pt] {#1};
  \draw[eje] (0,0) -- (0,#4) node[left] {#2};
  \node[below left] at (0,0) {$0$};}
% \proyeccion{(x,y)}{etq x}{etq y}: líneas guía a los ejes con etiquetas
\newcommand{\proyeccion}[3]{%
  \draw[guia] let \p1=#1 in (\x1,0) node[below] {#2} |- (0,\y1) node[left] {#3};}

%   \begin{document}
%   \begin{tikzpicture}[tcee]
%     \ejes{Q}{P}{6}{5}
%     \draw[curva1] (0.5,4.5) -- (5,0.8) node[etiqueta,right] {$D$};
%     \capa{2}{\draw[curva2] ...;}   % en el vídeo aparece en el paso 2
%   \end{tikzpicture}
%   \end{document}
%
% Paleta: la de la web (morado, dorado, salvia) más tres colores de apoyo
% distinguibles en blanco y negro por el trazo.
% ==========================================================================
\usepackage[T1]{fontenc}
\usepackage{newpxtext,newpxmath}
\usepackage{xcolor}
\usepackage{tikz,pgfplots}
\pgfplotsset{compat=1.18}
\usetikzlibrary{arrows.meta,calc,positioning,patterns,decorations.pathreplacing,intersections,shapes.misc,backgrounds}

\definecolor{tceemorado}{HTML}{5F2987}
\definecolor{tceedorado}{HTML}{B8860B}
\definecolor{tceeverde}{HTML}{3C7D3A}
\definecolor{tceeazul}{HTML}{1F5F99}
\definecolor{tceerojo}{HTML}{B22222}
\definecolor{tceegris}{HTML}{767171}
\definecolor{tceesalvia}{HTML}{E2EFD9}

% Capas para el vídeo: \capa{n}{…} dibuja su contenido solo a partir del paso n.
% En el PDF, la web y Word se ve todo; generar-video.py compila el gráfico una
% vez por paso con \def\tceepaso{k} para que cada elemento aparezca al narrarlo.
\providecommand{\tceepaso}{1000}
\newcommand{\capa}[2]{\ifnum#1>\tceepaso\relax\else#2\fi}

\tikzset{
  tcee/.style={line cap=round,line join=round,font=\small,>={Stealth[length=2.2mm]}},
  eje/.style={->,thick,black},
  curva1/.style={very thick,tceemorado},
  curva2/.style={very thick,tceeazul},
  curva3/.style={very thick,tceeverde},
  curva4/.style={very thick,tceedorado},
  curvanueva/.style={very thick,tceerojo},
  desplazada/.style={very thick,dashed},
  guia/.style={thin,dashed,tceegris},
  punto/.style={circle,fill=black,inner sep=1.3pt},
  etiqueta/.style={font=\small,inner sep=2pt},
  flecha/.style={->,thick,tceerojo},
  area/.style={fill=tceemorado,fill opacity=0.18,draw=none},
  area2/.style={fill=tceedorado,fill opacity=0.22,draw=none},
}

% \ejes{etiqueta x}{etiqueta y}{largo x}{alto y}
\newcommand{\ejes}[4]{%
  \draw[eje] (0,0) -- (#3,0) node[below left=2pt and 0pt] {#1};
  \draw[eje] (0,0) -- (0,#4) node[left] {#2};
  \node[below left] at (0,0) {$0$};}
% \proyeccion{(x,y)}{etq x}{etq y}: líneas guía a los ejes con etiquetas
\newcommand{\proyeccion}[3]{%
  \draw[guia] let \p1=#1 in (\x1,0) node[below] {#2} |- (0,\y1) node[left] {#3};}

%   \begin{document}
%   \begin{tikzpicture}[tcee]
%     \ejes{Q}{P}{6}{5}
%     \draw[curva1] (0.5,4.5) -- (5,0.8) node[etiqueta,right] {$D$};
%     \capa{2}{\draw[curva2] ...;}   % en el vídeo aparece en el paso 2
%   \end{tikzpicture}
%   \end{document}
%
% Paleta: la de la web (morado, dorado, salvia) más tres colores de apoyo
% distinguibles en blanco y negro por el trazo.
% ==========================================================================
\usepackage[T1]{fontenc}
\usepackage{newpxtext,newpxmath}
\usepackage{xcolor}
\usepackage{tikz,pgfplots}
\pgfplotsset{compat=1.18}
\usetikzlibrary{arrows.meta,calc,positioning,patterns,decorations.pathreplacing,intersections,shapes.misc,backgrounds}

\definecolor{tceemorado}{HTML}{5F2987}
\definecolor{tceedorado}{HTML}{B8860B}
\definecolor{tceeverde}{HTML}{3C7D3A}
\definecolor{tceeazul}{HTML}{1F5F99}
\definecolor{tceerojo}{HTML}{B22222}
\definecolor{tceegris}{HTML}{767171}
\definecolor{tceesalvia}{HTML}{E2EFD9}

% Capas para el vídeo: \capa{n}{…} dibuja su contenido solo a partir del paso n.
% En el PDF, la web y Word se ve todo; generar-video.py compila el gráfico una
% vez por paso con \def\tceepaso{k} para que cada elemento aparezca al narrarlo.
\providecommand{\tceepaso}{1000}
\newcommand{\capa}[2]{\ifnum#1>\tceepaso\relax\else#2\fi}

\tikzset{
  tcee/.style={line cap=round,line join=round,font=\small,>={Stealth[length=2.2mm]}},
  eje/.style={->,thick,black},
  curva1/.style={very thick,tceemorado},
  curva2/.style={very thick,tceeazul},
  curva3/.style={very thick,tceeverde},
  curva4/.style={very thick,tceedorado},
  curvanueva/.style={very thick,tceerojo},
  desplazada/.style={very thick,dashed},
  guia/.style={thin,dashed,tceegris},
  punto/.style={circle,fill=black,inner sep=1.3pt},
  etiqueta/.style={font=\small,inner sep=2pt},
  flecha/.style={->,thick,tceerojo},
  area/.style={fill=tceemorado,fill opacity=0.18,draw=none},
  area2/.style={fill=tceedorado,fill opacity=0.22,draw=none},
}

% \ejes{etiqueta x}{etiqueta y}{largo x}{alto y}
\newcommand{\ejes}[4]{%
  \draw[eje] (0,0) -- (#3,0) node[below left=2pt and 0pt] {#1};
  \draw[eje] (0,0) -- (0,#4) node[left] {#2};
  \node[below left] at (0,0) {$0$};}
% \proyeccion{(x,y)}{etq x}{etq y}: líneas guía a los ejes con etiquetas
\newcommand{\proyeccion}[3]{%
  \draw[guia] let \p1=#1 in (\x1,0) node[below] {#2} |- (0,\y1) node[left] {#3};}

\begin{document}
\begin{tikzpicture}[tcee,
    cab/.style={font=\small\bfseries,align=center,text width=3.1cm},
    celda/.style={font=\small,align=center,text width=3.1cm},
    rejilla/.style={thick,tceegris}]
  % 1: ecuación y cabecera
  \capa{1}{\fill[tceesalvia] (0,0) rectangle (10.2,1.4);}
  \capa{1}{\node[font=\normalsize] at (5.1,0.75)
    {$\underbrace{\partial x_i/\partial p_i}_{\text{ET}}=\underbrace{\partial h_i/\partial p_i}_{\text{ES}}\ \underbrace{-\,x_i\,\partial x_i/\partial\overline{W}}_{\text{ER}}$};}
  \capa{1}{\draw[rejilla] (0,0) rectangle (10.2,1.4);}
  \capa{1}{\node[cab] at (1.7,-0.4) {Efecto sustitución};}
  \capa{1}{\node[cab] at (5.1,-0.4) {Efecto renta};}
  \capa{1}{\node[cab] at (8.5,-0.4) {Efecto total};}
  \capa{1}{\draw[rejilla] (0,-0.8) -- (10.2,-0.8);}
  % 2: bien normal
  \capa{2}{\node[celda] at (1.7,-1.3) {${\leq}\,0$};}
  \capa{2}{\node[celda,text=tceeverde] at (5.1,-1.3) {${<}\,0$: bien normal};}
  \capa{2}{\node[celda,text=tceeazul] at (8.5,-1.3) {${<}\,0$: ordinario};}
  \capa{2}{\draw[rejilla] (0,-1.8) -- (10.2,-1.8);}
  % 3: bien inferior ordinario (domina el efecto sustitución)
  \capa{3}{\node[celda] at (1.7,-2.3) {${\leq}\,0$};}
  \capa{3}{\node[celda,text=tceedorado] at (5.1,-2.3) {${>}\,0$: bien inferior};}
  \capa{3}{\node[celda,text=tceeazul] at (8.5,-2.3) {${\leq}\,0$: ordinario\\ \footnotesize si $|\text{ES}|\geq\text{ER}$};}
  \capa{3}{\draw[rejilla] (0,-2.8) -- (10.2,-2.8);}
  % 4: bien Giffen (domina el efecto renta)
  \capa{4}{\node[celda] at (1.7,-3.3) {${\leq}\,0$};}
  \capa{4}{\node[celda,text=tceedorado] at (5.1,-3.3) {${>}\,0$: bien inferior};}
  \capa{4}{\node[celda,text=tceerojo] at (8.5,-3.3) {${>}\,0$: \textsc{Giffen}\\ \footnotesize si $\text{ER}>|\text{ES}|$};}
  \capa{4}{\draw[rejilla] (0,-3.8) -- (10.2,-3.8);}
  % rejilla vertical, que crece con las filas
  \foreach \k/\a/\b in {1/0/-0.8,2/-0.8/-1.8,3/-1.8/-2.8,4/-2.8/-3.8}{
    \capa{\k}{\foreach \x in {0,3.4,6.8,10.2} \draw[rejilla] (\x,\a) -- (\x,\b);}
  }
\end{tikzpicture}
\end{document}
